\documentclass[a4paper,10pt]{article}
\usepackage[utf8]{inputenc}
\usepackage[T1]{fontenc}
\usepackage{lmodern}
\usepackage[margin=1.4cm, top=1.2cm, bottom=1.2cm]{geometry}
\usepackage{hyperref}
\usepackage{xcolor}
\usepackage{enumitem}
\usepackage{graphicx}
\usepackage{titlesec}

% Minimal clean styling
\hypersetup{colorlinks=true, urlcolor=blue!70!black, linkcolor=blue!70!black}
\pagestyle{empty}
\setlength{\parindent}{0pt}

% Section formatting
\titleformat{\section}{\large\bfseries\color{black}}{}{0em}{}[\titlerule]
\titlespacing*{\section}{0pt}{6pt}{4pt}

% Colors
\definecolor{accent}{RGB}{20, 50, 120}

% Role command
\newcommand{\role}[4]{%
  \vspace{2pt}\noindent
  \begin{minipage}[t]{0.70\textwidth}\raggedright\textbf{#1} \textbullet\ \textit{#3}\end{minipage}%
  \hfill
  \begin{minipage}[t]{0.28\textwidth}\raggedleft\footnotesize\color{gray}#2 \textbullet\ #4\end{minipage}%
  \vspace{1pt}%
}
\setlist[itemize]{leftmargin=1.2em, itemsep=1.5pt, topsep=2pt, parsep=0pt}

\begin{document}

% ================= HEADER =================
\begin{minipage}[c]{0.80\textwidth}
  {\Huge\bfseries Prabhavit Gupta}\\[3pt]
  {\color{accent}\large AI Engineer $\&$ Generative AI Engineer $\vert$ LLM, RAG, Vector DB, MLOps $\&$ Python}\\[6pt]
  {\footnotesize +39 351 610 7807 ~$\vert$~
  \href{mailto:prabhavitg@gmail.com}{prabhavitg@gmail.com} ~$\vert$~
  \href{https://www.linkedin.com/in/prabhavit/}{linkedin.com/in/prabhavit} ~$\vert$~
  \href{https://github.com/prab-gupta}{github.com/prab-gupta} ~$\vert$~
  Venezia, Italia (Residente Permanente UE $\vert$ Disponibilità Immediata)}
\end{minipage}\hfill

% ================= SUMMARY =================
\section{Summary}
AI Engineer e Generative AI Engineer con 3 anni di esperienza commerciale nello sviluppo, testing e industrializzazione di soluzioni end-to-end basate su intelligenza artificiale, Large Language Models (LLM) e Machine Learning. Esperto nella progettazione di architetture RAG (Retrieval-Augmented Generation) avanzate con Vector Database (Qdrant, Pinecone, pgvector), tecniche di Prompt Engineering e generazione di dati sintetici (Synthetic Data Generation) per benchmark e validazione affidabile dei sistemi AI. Comprovata esperienza nella collaborazione con team di Data e Platform Engineering per il deployment in produzione di microservizi containerizzati (Docker, Kubernetes) su cloud AWS/Azure/GCP e nell'implementazione di pratiche MLOps. Residente permanente a Venezia, con ottimo inglese professionale (C1) e italiano lavorativo.

% ================= EXPERIENCE =================
\section{Experience}

\role{\href{https://flowspace.co}{FlowSpace.co}}{Mar 2024 - Aug 2026}{Lead AI Engineer $\&$ GenAI Architect (Freelance)}{Londra, UK (remoto)}
\begin{itemize}
  \item Progettato e rilasciato in produzione architetture generative e pipeline RAG per oltre 4.500 utenti enterprise attivi, integrando LLM commerciali e open-source con Vector Database (Qdrant, pgvector).
  \item Implementato tecniche avanzate di Prompt Engineering, structured outputs e catene di reasoning per eliminare allucinazioni e garantire latenze <1.8s in ambienti di produzione ad alto volume.
  \item Sviluppato pipeline di Synthetic Data Generation per generare dataset di test e valutare oggettivamente le metriche di retrieval (hit rate, MRR) e generazione (faithfulness, answer relevance).
  \item Coordinato con i team Data Engineering e Platform Engineering l'orchestrazione di microservizi FastAPI su Docker e AWS/GCP, integrando flussi di CI/CD e monitoraggio MLOps.
  \item Collaborato con stakeholder tecnici e di business per tradurre requisiti di dominio in specifiche architetturali e linee guida di governance AI.
  \item \textit{Tecnologie: Python, Generative AI, PyTorch, RAG, Qdrant, Pinecone, FastAPI, Docker, AWS, Databricks, MLOps, Git}
\end{itemize}

\role{Urlink}{Apr 2025 - Jul 2025}{AI $\&$ Machine Learning Engineer (Freelance)}{Barcellona, Spagna (remoto)}
\begin{itemize}
  \item Sviluppato modelli e pipeline per l'estrazione semantica, classificazione automatica e normalizzazione di dati testuali non strutturati.
  \item Realizzato pipeline di test automatizzati con dati sintetici per verificare la robustezza dei modelli contro edge cases ed errori di ingestion.
  \item Progettato API REST e microservizi per servire inferenze ML in tempo reale con validazione rigorosa degli schemi tramite Pydantic.
  \item \textit{Tecnologie: Python, Scikit-learn, PyTorch, Pandas, PostgreSQL, Docker, REST APIs, Git}
\end{itemize}

\role{Consulente Tecnico Indipendente}{Sep 2023 - Mar 2024}{AI Solutions Engineer}{Venezia, Italia}
\begin{itemize}
  \item Progettato e sviluppato soluzioni personalizzate di Generative AI, agenti conversazionali e workflow RAG per PMI europee.
  \item Guidato le fasi di prototipazione rapida, testing su benchmark dedicati e deployment su infrastrutture cloud sicure.
  \item Redatto documentazione tecnica completa, linee guida di prompt engineering e procedure operative per il monitoraggio dei modelli.
  \item \textit{Tecnologie: Python, OpenAI API, LangChain, LlamaIndex, Vector DB, FastAPI, Linux, Git}
\end{itemize}

\role{\href{https://tractionboard.co}{Tractionboard.co}}{Aug 2023 - Oct 2023}{Full-Stack Automation Developer (Stage)}{Utrecht, Paesi Bassi (remoto)}
\begin{itemize}
  \item Sviluppato connettori webhook e pipeline di ingestion dati per piattaforme SaaS di analytics.
  \item \textit{Tecnologie: Python, JavaScript, REST APIs, Git}
\end{itemize}

% ================= EDUCATION =================
\section{Education}
\role{Università Ca' Foscari di Venezia}{Ott 2020 - Giu 2024}{Laurea Triennale in Filosofia, Studi Internazionali ed Economia, Minor in Imprenditorialità}{Venezia, Italia}

% ================= CERTIFICATIONS =================
\section{Certifications}
Microsoft: Introduction to building AI agents \textbullet\ Microsoft Copilot Studio Agent Academy \textbullet\ Microsoft Power Automate \textbullet\ Tray.AI Build Practitioner (Sep 2024) \textbullet\ Tray.AI Build Foundations

% ================= ADDITIONAL INFORMATION =================
\section{Additional Information}
\begin{itemize}
  \item \textbf{Inquadramento:} Residente permanente UE a Venezia, Italia. Disponibile per contratti CCNL o collaborazione con Partita IVA.
  \item \textbf{Lingue:} Italiano (lavorativo), Inglese (fluente C1 professionale), Hindi (madrelingua).
  \item \textbf{Disponibilità:} Immediata (0 giorni di preavviso) per posizioni full-time remote / ibride.
\end{itemize}

% ================= CORE SKILLS =================
\section{Core Skills}
\begin{itemize}
  \item \textbf{Generative AI $\&$ LLM:} Large Language Models, Prompt Engineering, Synthetic Data Generation, Structured Outputs, Evals.
  \item \textbf{RAG $\&$ Vector Databases:} Retrieval-Augmented Generation, Qdrant, Pinecone, pgvector, Chunking Semantico, Hybrid Search.
  \item \textbf{Machine Learning $\&$ Linguaggi:} Python 3.12, PyTorch, Scikit-learn, TensorFlow, Databricks, FastAPI, Pydantic, SQL.
  \item \textbf{MLOps, Cloud $\&$ Architetture:} Docker, Kubernetes, AWS, Microsoft Azure, Google Cloud Platform, Microservizi, CI/CD, Git.
  \item \textbf{Metodologie:} Agile/Scrum, Collaboration Cross-funzionale, Redazione Architetturale, Problem Solving Analitico.
\end{itemize}

\end{document}
